\documentclass[a4paper,12pt]{amsbook}
\usepackage{physics}
\usepackage{amsthm}
\usepackage{amsmath,amssymb,amsthm}
\usepackage{mathtools}
\usepackage{mathpazo}
\usepackage{microtype}
\providecommand{\tb}{\textbackslash}
\newcommand{\0}{{$|0\rangle$}}
\newcommand{\1}{{$|1\rangle$}}
\usepackage{graphicx}
\usepackage{amsfonts}
\usepackage[utf8]{inputenc}
\usepackage[T1]{fontenc}
\usepackage{geometry}
\geometry{a4paper, margin=1in}
\usepackage{listings}
\usepackage{xcolor}
\pagecolor{black}
\color{white}
\usepackage{parskip}
\setlength{\parindent}{0pt}
\newtheorem{lemma}{Lemma}

\begin{document}
\author{Vivek Soorya Maadoori}
\title{Density Operators cant have non-zero elements in all diagonal positions}
\maketitle
\begin{lemma}
A density operator of a pure state cannot have non-zero elements in all diagonal positions.
\end{lemma}
\begin{proof}

    \footnote{Density operators of a pure state are projectors}
    Density operators are positive operators and so are diagonalizable. 
    Assume the diagonalized density matrix $\rho_D$. Assume one nonzero entry. 

    Let $\rho_D = |\psi\rangle\langle\psi| \text{ for some } |\psi\rangle $.
    For the state $|\psi\rangle$, let the entries be $x_n, n \in \mathcal{N}$.
    For any indices i, j such that i $\neq$ j, either $x_i$ is 0 or $x_j$ is zero, since all off diagonal elements are zero. Thus, at least one of $x_i$ and $x_j$ are zero for any distinct pair of indices. 
    $\Rightarrow$ There is only one index i such that $x_i \neq 0$
    $\Rightarrow$ There is only one positive diagonal element in $\rho_D$

    Thus, given the non diagonal elements of a projector are 0 and that the projector is a positive matrix, There can only be one non-zero diagonal element. 
\end{proof}
\end{document}
